\chapter{Numerical Reconstruction}

\section*{Motivation}

\begin{remark}
In the curvature-based framework, geometric objects are not given
explicitly as coordinate sets. Instead, they are reconstructed from
curvature information by integration.
\end{remark}

\begin{remark}
As a consequence, numerical methods are not merely implementation
details, but an essential part of the modelling process.
\end{remark}

---

\section{Reconstruction Principle}

\begin{definition}[Reconstruction from curvature]
Let $\kappaf(s)$ be a curvature function and let
\[
(x_0,y_0,\theta_0)
\]
be a start pose. The corresponding curve is obtained by solving
\[
\theta'(s) = \kappaf(s),
\qquad
\gamma'(s) = (\cos\theta(s),\sin\theta(s)).
\]
\end{definition}

\begin{remark}
Even if $\kappaf$ is given analytically, closed-form solutions for
$\gamma$ are generally not available.
\end{remark}

---

\section{Discretization}

\begin{definition}[Discretization]
A discretization of the parameter interval $[0,L]$ is a finite sequence
\[
0 = s_0 < s_1 < \dots < s_n = L.
\]
\end{definition}

\begin{definition}[Discrete reconstruction]
A discrete reconstruction computes approximate values
\[
\gamma(s_i), \quad \theta(s_i)
\]
at the discretization points.
\end{definition}

\begin{remark}
The choice of discretization strongly influences both computational cost
and geometric accuracy.
\end{remark}

\TODO{Formal criteria for discretization strategies (uniform vs adaptive).}

---

\section{Integration Methods}

\begin{remark}
The reconstruction problem can be formulated as an initial value problem
for a system of ordinary differential equations.
\end{remark}

\subsection{Basic Scheme}

\begin{definition}[Explicit step]
Given $(\gamma_i,\theta_i)$ at $s_i$, an explicit step computes
\[
\theta_{i+1} = \theta_i + \kappaf(s_i)\,\Delta s,
\]
\[
\gamma_{i+1} = \gamma_i + (\cos\theta_i,\sin\theta_i)\,\Delta s.
\]
\end{definition}

\begin{remark}
This corresponds to a first-order method and is generally insufficient
for high-precision reconstruction.
\end{remark}

\subsection{Higher-Order Methods}

\TODO{Runge--Kutta methods for coupled systems.}

\TODO{Romberg integration for curvature integration.}

\TODO{Adaptive step size control.}

\OPEN{Separation of curvature integration and geometric integration.}

---

\section{Error Analysis}

\begin{definition}[Local error]
The local error of a step is the deviation introduced at a single
integration step.
\end{definition}

\begin{definition}[Global error]
The global error is the accumulated deviation over the entire interval.
\end{definition}

\begin{remark}
Errors arise from:
\begin{itemize}
\item discretization,
\item numerical integration,
\item approximation of curvature laws.
\end{itemize}
\end{remark}

\RESEARCH{Quantitative bounds for error propagation in curvature-driven reconstruction.}

\section{Sensitivity}

\begin{remark}
The reconstruction process is sensitive to perturbations in curvature.
Small deviations in $\kappaf$ may lead to significant geometric drift.
\end{remark}

\RESEARCH{Sensitivity analysis with respect to curvature perturbations.}

\OPEN{Relation between curvature noise and positional uncertainty.}

---

\section{Sampling and Representation}

\begin{definition}[Sampling]
Sampling refers to the extraction of a finite set of points from a
continuous curve representation.
\end{definition}

\begin{remark}
Sampling is required for:
\begin{itemize}
\item visualization,
\item collision detection,
\item export to point-based formats.
\end{itemize}
\end{remark}

\TODO{Error-controlled sampling strategies.}

\OPEN{Difference between sampling for visualization and sampling for computation.}

---

\section{Reconstruction in the Sparse Model}

\begin{remark}
In the sparse alignment model, reconstruction is performed element-wise.
Each element contributes a local curvature law, which is integrated
sequentially.
\end{remark}

\begin{definition}[Sequential reconstruction]
Let $(e_1,\dots,e_n)$ be a sparse alignment. The reconstruction proceeds
by integrating each element in sequence, using the end state of the
previous element as initial state for the next.
\end{definition}

\OPEN{Handling of numerical drift across element boundaries.}

---

\section{Numerical Stability}

\begin{remark}
Numerical stability is a central concern, especially for long alignments
or highly curved geometries.
\end{remark}

\RESEARCH{Stability criteria for curvature-based integration.}

\TODO{Techniques for drift correction and re-normalization.}

---

\section{Connection to Optimization}

\begin{remark}
Numerical reconstruction is closely linked to optimization problems.
In fitting and design tasks, reconstruction is performed repeatedly
within iterative algorithms.
\end{remark}

\TODO{SQP formulation for alignment fitting.}

\RESEARCH{Connection to optimal control problems.}

---

\section{Summary}

\begin{summary}
Numerical reconstruction is the process of deriving geometric
representations from curvature data by integration.

It introduces discretization, approximation, and error propagation as
intrinsic components of the modelling framework.

Therefore, numerical methods are not merely technical tools, but an
essential layer connecting mathematical definitions with practical
geometry.
\end{summary}

\section{Numerical Reconstruction from Curvature}

\subsection{Problem Statement}

\begin{remark}
Given a curvature function
\[
\kappa : [0,L] \to \R,
\]
and initial conditions
\[
\gamma(0), \quad \theta(0),
\]
the geometric alignment is obtained by solving:
\[
\gamma'(s) = (\cos\theta(s), \sin\theta(s)),
\qquad
\theta'(s) = \kappa(s).
\]
\end{remark}

\begin{remark}
This defines a nonlinear initial value problem.
\end{remark}

---

\subsection{Discrete Approximation}

\begin{definition}[Discretization]
Let
\[
0 = s_0 < s_1 < \dots < s_N = L
\]
be a partition with step sizes \(\Delta s_i\).
\end{definition}

\begin{remark}
We approximate:
\[
\kappa_i \approx \kappa(s_i).
\]
\end{remark}

---

\subsection{Forward Integration Scheme}

\begin{definition}[Explicit integration]
Given \((\gamma_i, \theta_i)\), define:
\[
\theta_{i+1} = \theta_i + \kappa_i \Delta s_i,
\]
\[
\gamma_{i+1} = \gamma_i + \Delta s_i
(\cos\theta_i, \sin\theta_i).
\]
\end{definition}

\begin{remark}
This corresponds to an explicit Euler scheme.
\end{remark}

---

\subsection{Midpoint Scheme}

\begin{definition}[Midpoint method]
Define:
\[
\theta_{i+\frac{1}{2}} = \theta_i + \frac{1}{2}\kappa_i \Delta s_i,
\]
\[
\gamma_{i+1} = \gamma_i + \Delta s_i
(\cos\theta_{i+\frac{1}{2}}, \sin\theta_{i+\frac{1}{2}}),
\]
\[
\theta_{i+1} = \theta_i + \kappa_i \Delta s_i.
\]
\end{definition}

\begin{remark}
This improves accuracy and stability compared to the explicit scheme.
\end{remark}

---

\subsection{Extension to Pose3}

\begin{remark}
Including cant, we obtain:
\[
\phi_{i+1} = \phi_i + \omega_i \Delta s_i.
\]
\end{remark}

\begin{remark}
Thus, the full discrete system is:
\[
(\gamma_i, \theta_i, \phi_i)
\;\longrightarrow\;
(\gamma_{i+1}, \theta_{i+1}, \phi_{i+1}).
\]
\end{remark}

---

\subsection{Numerical Stability}

\begin{remark}
The reconstruction is sensitive to:
\begin{itemize}
\item step size \(\Delta s\),
\item noise in \(\kappa_i\),
\item and large coordinate values.
\end{itemize}
\end{remark}

\begin{remark}
To ensure stability:
\begin{itemize}
\item sufficiently small step sizes must be used,
\item curvature should be smoothed,
\item and local coordinate frames are recommended.
\end{itemize}
\end{remark}

---

\subsection{Floating Origin Technique}

\begin{remark}
For large coordinates, numerical precision degrades.
\end{remark}

\begin{definition}[Floating origin]
A local coordinate system is introduced such that:
\[
\gamma_i^{\text{local}} = \gamma_i - \gamma_{\text{ref}}.
\]
\end{definition}

\begin{remark}
This reduces numerical errors in computation and visualization.
\end{remark}

---

\subsection{Consistency with CRS}

\begin{remark}
Reconstruction depends on:
\begin{itemize}
\item initial position,
\item initial orientation,
\item and coordinate system definition.
\end{itemize}
\end{remark}

\begin{remark}
Thus, numerical reconstruction must be performed within a consistent CRS.
\end{remark}

---

\subsection{Interpretation}

\begin{summary}
Numerical reconstruction provides the link between curvature-based
models and geometric representation.

It transforms abstract functions into concrete spatial geometry,
while preserving the structure required for analysis and optimization.
\end{summary}
